\section{Discussion}

\subsection{The Performance Variation Hierarchy}

The locked-configuration analysis yielded the hierarchy

\begin{center}
Data source heterogeneity $\gg$ Splitting method $\approx$ In-domain algorithm choice
\end{center}

with the data source heterogeneity gap ($\Delta = 0.390$) an order of magnitude larger than the split-method gap ($\Delta = 0.046$), 
itself comparable to the in-domain algorithm spread (MCC range = 0.063). The hierarchy was established on a single Ames benchmark; 
the LDO $\Delta$ also reflects asymmetric test sizes and class distributions and should be read as an upper bound; and algorithm robustness was uneven under source shift, 
so algorithm choice may matter much more in cross-source deployment than the in-domain comparison suggests.

Importantly, the LDO experiment is not a deployment scenario.
In practice, models are trained on all available data regardless of provenance, and the genotoxic potential of a chemical is fundamentally a property of its molecular structure,
not of which database it was registered in. The point of cross-source evaluation is rather diagnostic:
a large LDO drop reveals that apparent in-domain performance depends on the joint composition of training and test sets rather than on universally learned structure-activity relationships.
This matters because new chemicals registered under EU REACH or K-REACH may occupy regions of chemical space that are sparsely represented in pharmaceutical benchmarks,
and a benchmark MCC of 0.67 on a mixed-source test set provides little guarantee about the model's reliability on such compounds.
Cross-source LDO acts as a stress test that exposes this dependency and gives a more honest lower bound on out-of-distribution performance.

\subsection{What Drives the Data Source Heterogeneity Gap?}

Four non-exclusive hypotheses can be identified, in roughly decreasing order of parsimony.
\textbf{Prior shift} from the 19.5-fold prevalence gap can mechanically lower MCC by 
miscalibrating the default threshold, consistent with XGB reaching sensitivity 0.062 in the 
industrial$\to$drug direction. \textbf{Benchmark enrichment} is the most direct explanation 
for the prevalence gap itself: published Ames benchmarks such as 
Hansen\cite{Hansen_2009} were assembled to include a balanced set of known mutagens and 
non-mutagens for evaluation purposes, whereas industrial REACH/K-REACH dossiers reflect the 
prevalence of mutagenicity in the broader chemical economy, which is genuinely low. 
\textbf{Data quality and labeling heterogeneity}---differing protocols, strain panels, S9 
mix conditions, and curation conventions between drug and industrial sources---can shift 
both prevalence and decision boundaries. \textbf{Chemical space difference} would mean that 
even a perfectly calibrated cross-source model generalizes poorly because the two sources 
occupy systematically different regions of descriptor space; published QSAR-curation work 
documents the broad consequences of such differences for model 
transportability.\cite{Tropsha_2010,Fourches_2010}

% TODO[smyoon]: References Lipinski_2004 and Veber_2002 (originally cited here for 
% "chemical space difference") are about drug-likeness/oral bioavailability, not QSAR 
% generalization across chemical-space partitions. They have been removed; if you want a 
% chemistry-space-specific citation here, candidates include Sheridan 2013 (already cited 
% as Sheridan_2013), Tropsha 2010 (Tropsha_2010), or a recent applicability-domain review. 
% Confirm before resubmission.

To estimate the contribution of prior shift, each LDO direction was evaluated at three 
thresholds: \textit{default} ($t = 0.5$), \textit{prior-corrected} ($t$ = training positive 
rate, optimal under pure prior shift), and \textit{oracle} ($t$ maximizing MCC on the test 
set, an upper bound on threshold-based recalibration). Two complementary analyses are 
provided: Table~S7 uses Morgan FP-1024 + 13 physicochemical descriptors on the full LDO 
partition, and Table~S8 uses the locked-configuration 138-feature compact set on a reduced 
cross-source partition (drug train $n=5{,}027$; industrial train $n=5{,}821$).

The headline quantitative result comes from Table~S8: across all six 
algorithm-by-direction combinations on the locked-configuration features, threshold tuning 
explains 30.9--35.3\,\% of the total LDO gap (default-threshold MCC vs.\ in-domain scaffold-CV 
MCC of 0.6241), leaving a residual of $\Delta_{\mathrm{residual}} = 0.15$--$0.37$ MCC that 
cannot be recovered by recalibration. The drug $\to$ industrial direction has a larger 
residual ($\Delta_{\mathrm{residual}} = 0.34$--$0.37$) than industrial $\to$ drug 
($\Delta_{\mathrm{residual}} = 0.15$--$0.18$), consistent with industrial-to-drug 
generalization being supported by a larger and more diverse training set. The Table~S7 
analysis on the full partition with FP-1024 features shows the same qualitative pattern 
(prior-shift correction yields modest improvements, oracle MCC plateaus well below in-domain 
performance), confirming that the gap-decomposition conclusion is robust to feature choice.

In summary, prior shift accounts for roughly one-third of the LDO collapse, and the 
remaining two-thirds reflect genuine source shift---some combination of chemical space 
difference, data quality and labeling heterogeneity, and benchmark enrichment---that 
threshold tuning alone cannot eliminate. Published Ames QSAR performance is therefore highly 
sensitive to source composition, and threshold tuning is not a sufficient remedy.

\textit{A note on the RF robustness finding.} The default-threshold MCC values for RF 
differ across analyses: Table~\ref{tab:ldo} (full partition, compact features) shows RF 
industrial $\to$ drug = 0.325, Table~S7 (full partition, FP-1024 features) shows 0.102, and 
Table~S8 (reduced partition, compact features) shows 0.350. Tables~\ref{tab:ldo} and S8 are 
broadly consistent (difference $\leq 0.03$), supporting the main-text claim that RF retains 
useful cross-source classification ability in the harder direction when run with 
domain-knowledge features. The Table~S7 result on FP-1024 features differs more 
substantially, suggesting that the structural-alert component of the compact set carries 
non-trivial cross-source information that fingerprints do not fully replicate. The 
algorithm-comparison pattern at default threshold should therefore be read as 
feature-conditional. The within-row gap decomposition in Table~S8 does not depend on 
algorithm choice in this way: all three algorithms show the same 30.9--35.3\,\% prior-shift 
share.

% TODO[smyoon]: Table S8 (locked-config compact features, reduced partition) addresses the 
% earlier "re-run threshold decomposition with 138-feature compact" TODO. The remaining 
% open question is reconciling the reduced N (5027/5821 in S8) with the full N (6416/7144 
% in Table 5). If both can be regenerated on a single shared partition, the cleanest 
% revision is to do so and merge S7/S8 into one table; if not, the methodological note in 
% S8 documenting the partition difference is the minimum acceptable disclosure.

\subsection{In-Vivo Micronucleus and Tiered Testing}

The dissociation between ranking ability (ROC-AUC = 0.860, PR-AUC = 0.237) and classification 
reliability (MCC CI crossing zero) means that, on the present test set, in-vivo MN QSAR with 
2D descriptors does not provide stable binary classification at any fixed threshold; a wider 
external evaluation would be needed to determine whether this reflects an intrinsic limit of 
2D representations or a sample-size limit (only 12 positives in the test set). The 8.5-fold 
PR-AUC enrichment over baseline suggests value as one input in tiered screening under ICH 
M7\cite{ICH_M7_2017} or REACH,\cite{EU_REACH_2006,Dobo_2012} provided downstream decisions 
do not rely on threshold-based classification alone. The strong preprocessing sensitivity 
($+0.22$ MCC from salt stripping) means reported in-vivo performance depends critically on 
preprocessing choices that are often undocumented.\cite{Fourches_2016} Mechanistically, 
micronucleus induction depends on ADME, metabolic activation, tissue distribution, and DNA 
repair,\cite{Tweats_2007,Kirkland_2011} none of which 2D descriptors capture.

\subsection{Cross-Endpoint Overlap and Preprocessing}

The high compound overlap between endpoints ($>$90\% shared with Ames) implies that performance differences (Ames MCC = 0.555 vs.\ in-vivo 0.311) reflect endpoint predictability from 2D structure rather than 
differing chemical spaces.\cite{Tropsha_2010} Multi-task learning may be a productive direction.\cite{Ramsundar_2015} 
Salt stripping was beneficial for in-vivo, mixed for Ames, and harmful for in-vitro---possibly because some counterions are themselves clastogenically relevant. 
The QSAR curation literature\cite{Fourches_2010,Fourches_2016} emphasizes standardization, but the present results suggest the optimal preprocessing strategy is endpoint-dependent.

\subsection{Practical Recommendations}

Based on these findings, we propose the following as useful additions to current reporting 
practices for genotoxicity QSAR (rather than as a fixed standard from a single study): 
(i)~scaffold-based splitting; 
(ii)~when source labels exist, cross-source LDO performance with disclosed source composition; 
(iii)~per-endpoint preprocessing optimization, with preprocessing flags reported alongside 
metrics; 
(iv)~bootstrap CIs and AD coverage; 
(v)~for class-imbalanced endpoints, both threshold-free (ROC-AUC, PR-AUC) and 
threshold-dependent (MCC, balanced accuracy) metrics; 
(vi)~testing algorithm superiority claims under source shift, not just in-domain---in our 
data, in-domain MCC range across seven algorithms was 0.063 while cross-source MCC range 
across three tree-based algorithms reached 0.203.

Several constraints bound these findings. A single scaffold split was used rather than 
repeated holdout, the GCN did not include edge features or 
pre-training,\cite{Hu_2020_pretrain} and 3D or QM descriptors, metabolite predictions, 
within-AD vs.\ outside-AD stratification for LDO, and multiple-comparison correction for the 
McNemar tests\cite{Demsar_2006} were not evaluated. The Klimisch quality filter was applied 
only to the industrial-sourced data, since the drug-sourced benchmarks do not carry 
equivalent annotations; this asymmetry is a residual confound for the cross-source 
comparison. Finally, the threshold-decomposition analysis was run with a different feature 
set than the main pipeline, so its absolute MCC values are not directly comparable to 
Table~\ref{tab:ldo} (see Table~S7 footnote and the ``note on the RF robustness finding'' 
above).
